\textbf{Phần 1. Trắc nghiệm lựa chọn}
\setcounterpageref{ex}{0}

\begin{ex} %[3D2V72]
Tính giá trị của biểu thức: $A = 15\,000 \times 4 + 48\,000 : 6$.
\choice
{$65\,000$}
{\True $68\,000$}
{$70\,000$}
{$66\,000$}
\loigiai{
$15\,000 \times 4 = 60\,000$; $48\,000 : 6 = 8000$; $60\,000 + 8000 = 68\,000$.
\textbf{Đáp án của câu hỏi là B}
}
\end{ex}

\begin{ex} %[3D2V72]
Một cửa hàng có $84\,000\text{ kg}$ gạo. Ngày đầu bán $\dfrac{1}{4}$ số gạo, ngày thứ hai bán $\dfrac{1}{3}$ số gạo còn lại. Sau hai ngày cửa hàng còn lại bao nhiêu ki-lô-gam gạo?
\choice
{$21\,000\text{ kg}$}
{$35\,000\text{ kg}$}
{$48\,000\text{ kg}$}
{\True $42\,000\text{ kg}$}
\loigiai{
Ngày 1 bán: $84\,000 : 4 = 21\,000\text{ kg}$. Còn lại: $84\,000 - 21\,000 = 63\,000\text{ kg}$. Ngày 2 bán: $63\,000 : 3 = 21\,000\text{ kg}$. Còn lại: $63\,000 - 21\,000 = 42\,000\text{ (kg)}$.
\textbf{Đáp án của câu hỏi là D}
}
\end{ex}

\begin{ex} %[3D2V72]
Tìm $x$, biết: $(x - 15\,000) : 3 = 25\,000$.
\choice
{\True $90\,000$}
{$85\,000$}
{$95\,000$}
{$100\,000$}
\loigiai{
$x - 15\,000 = 25\,000 \times 3 = 75\,000 \Rightarrow x = 75\,000 + 15\,000 = 90\,000$.
\textbf{Đáp án của câu hỏi là A}
}
\end{ex}

\begin{ex} %[3D2V72]
Tính nhanh giá trị: $16\,500 \times 3 + 16\,500 \times 2 - 16\,500$.
\choice
{$65\,000$}
{$60\,000$}
{\True $66\,000$}
{$70\,000$}
\loigiai{
$16\,500 \times (3 + 2 - 1) = 16\,500 \times 4 = 66\,000$.
\textbf{Đáp án của câu hỏi là C}
}
\end{ex}

\begin{ex} %[3D2V72]
Trong phép chia cho 8 có thương là $11\,250$ và số dư là số dư lớn nhất có thể có. Số bị chia là
\choice
{$90\,000$}
{\True $90\,007$}
{$90\,008$}
{$90\,006$}
\loigiai{
Số dư lớn nhất là 7. Số bị chia: $8 \times 11\,250 + 7 = 90\,000 + 7 = 90\,007$.
\textbf{Đáp án của câu hỏi là B}
}
\end{ex}

\begin{ex} %[3D2V72]
Một mảnh vườn hình chữ nhật có chu vi $64\,000\text{ m}$, chiều rộng bằng $\dfrac{1}{3}$ chiều dài. Chiều dài mảnh vườn là
\choice
{$8000\text{ m}$}
{$16\,000\text{ m}$}
{$32\,000\text{ m}$}
{\True $24\,000\text{ m}$}
\loigiai{
Nửa chu vi: $64\,000 : 2 = 32\,000\text{ m}$. Rộng: $32\,000 : (3 + 1) = 8000\text{ m}$. Dài: $8000 \times 3 = 24\,000\text{ (m)}$.
\textbf{Đáp án của câu hỏi là D}
}
\end{ex}

\begin{ex} %[3D2V72]
Tính giá trị của biểu thức: $100\,000 - 18\,500 \times 4$.
\choice
{\True $26\,000$}
{$24\,000$}
{$28\,000$}
{$25\,000$}
\loigiai{
$18\,500 \times 4 = 74\,000$; $100\,000 - 74\,000 = 26\,000$.
\textbf{Đáp án của câu hỏi là A}
}
\end{ex}

\begin{ex} %[3D2V72]
Một nhà máy có 5 phân xưởng, mỗi phân xưởng sản xuất được $16\,800$ sản phẩm. Nhà máy đã xuất kho $55\,000$ sản phẩm. Số sản phẩm còn lại là
\choice
{$28\,000$ sản phẩm}
{$30\,000$ sản phẩm}
{\True $29\,000$ sản phẩm}
{$31\,000$ sản phẩm}
\loigiai{
Tổng: $16\,800 \times 5 = 84\,000$. Còn lại: $84\,000 - 55\,000 = 29\,000$ sản phẩm.
\textbf{Đáp án của câu hỏi là C}
}
\end{ex}

\begin{ex} %[3D2V72]
Tìm số tự nhiên $y$, biết: $y \times 5 + 24\,500 = 99\,500$.
\choice
{$14\,000$}
{\True $15\,000$}
{$16\,000$}
{$18\,000$}
\loigiai{
$y \times 5 = 99\,500 - 24\,500 = 75\,000 \Rightarrow y = 75\,000 : 5 = 15\,000$.
\textbf{Đáp án của câu hỏi là B}
}
\end{ex}

\begin{ex} %[3D2V72]
Một hình vuông có chu vi bằng chu vi hình chữ nhật có chiều dài $25\,000\text{ m}$ và chiều rộng $15\,000\text{ m}$. Cạnh của hình vuông là
\choice
{$40\,000\text{ m}$}
{$18\,000\text{ m}$}
{$22\,000\text{ m}$}
{\True $20\,000\text{ m}$}
\loigiai{
Chu vi hình chữ nhật: $(25\,000 + 15\,000) \times 2 = 80\,000\text{ m}$. Cạnh hình vuông: $80\,000 : 4 = 20\,000\text{ (m)}$.
\textbf{Đáp án của câu hỏi là D}
}
\end{ex}

\begin{ex} %[3D2V72]
Bác Năm có $90\,000\text{ kg}$ thóc. Bác chia đều vào 6 kho, sau đó bán đi 2 kho. Số thóc bác Năm còn lại là
\choice
{\True $60\,000\text{ kg}$}
{$30\,000\text{ kg}$}
{$45\,000\text{ kg}$}
{$50\,000\text{ kg}$}
\loigiai{
Mỗi kho có: $90\,000 : 6 = 15\,000\text{ kg}$. Còn lại 4 kho: $15\,000 \times 4 = 60\,000\text{ (kg)}$.
\textbf{Đáp án của câu hỏi là A}
}
\end{ex}

\begin{ex} %[3D2V72]
Tổng hai số là 80 000, số lớn gấp 4 lần số bé. Số lớn là
\choice
{$16\,000$}
{$60\,000$}
{\True $64\,000$}
{$70\,000$}
\loigiai{
Tổng số phần: $4 + 1 = 5$ phần. Số bé: $80\,000 : 5 = 16\,000$. Số lớn: $16\,000 \times 4 = 64\,000$.
\textbf{Đáp án của câu hỏi là C}
}
\end{ex}

\textbf{Phần 2. Trắc nghiệm đúng sai}
\setcounterpageref{ex}{0}

\begin{ex} %[3D2V72]
Xét tính đúng sai của các biểu thức tính toán phạm vi 100 000:
\choiceTF
{\True Biểu thức $20\,000 \times 4 : 2$ có giá trị bằng $40\,000$.}
{$100\,000 : (2 \times 5) = 100\,000$.}
{\True Trong biểu thức có ngoặc, thực hiện phép tính trong ngoặc trước.}
{Một số nhân với 0 luôn bằng chính số đó.}
\loigiai{
\begin{itemize}
    \item Ý a) Đúng: $80\,000 : 2 = 40\,000$.
    \item Ý b) Sai: Sai vì $100\,000 : 10 = 10\,000$.
    \item Ý c) Đúng: Quy tắc cơ bản.
    \item Ý d) Sai: Sai vì nhân với 0 luôn bằng 0.
\end{itemize}
\begin{itemize}
    \item a) \textbf{Đúng}.
    \item b) \textbf{Sai}.
    \item c) \textbf{Đúng}.
    \item d) \textbf{Sai}.
\end{itemize}
}
\end{ex}

\begin{ex} %[3D2V72]
Một kho hàng có $96\,000\text{ kg}$ phân bón chia đều vào 8 kho nhỏ:
\choiceTF
{\True Mỗi kho nhỏ chứa $12\,000\text{ kg}$ phân bón.}
{\True Ba kho nhỏ chứa tất cả $36\,000\text{ kg}$ phân bón.}
{Năm kho nhỏ chứa nhiều hơn 70 000 kg phân bón.}
{\True Nếu xuất đi 4 kho nhỏ thì trong kho còn lại một nửa số phân bón ban đầu.}
\loigiai{
\begin{itemize}
    \item Ý a) Đúng: $96\,000 : 8 = 12\,000\text{ kg}$.
    \item Ý b) Đúng: $12\,000 \times 3 = 36\,000\text{ kg}$.
    \item Ý c) Sai: Sai vì $12\,000 \times 5 = 60\,000 < 70\,000$ kg.
    \item Ý d) Đúng: $4 : 8 = \dfrac{1}{2}$.
\end{itemize}
\begin{itemize}
    \item a) \textbf{Đúng}.
    \item b) \textbf{Đúng}.
    \item c) \textbf{Sai}.
    \item d) \textbf{Đúng}.
\end{itemize}
}
\end{ex}

\textbf{Phần 3. Trả lời ngắn}
\setcounterpageref{ex}{0}

\begin{ex} %[3D2V72]
Tính giá trị của biểu thức: $18\,000 \times 4 - 32\,000 : 4$.
\shortans{64000}
\loigiai{
\textbf{Đáp số:} $64000$\\
$72\,000 - 8000 = 64\,000$.
}
\end{ex}

\begin{ex} %[3D2V72]
Một hình vuông có cạnh dài $16\,500\text{ m}$. Tính chu vi hình vuông theo mét.
\shortans{66000}
\loigiai{
\textbf{Đáp số:} $66000$\\
$16\,500 \times 4 = 66\,000\text{ (m)}$.
}
\end{ex}

\begin{ex} %[3D2V72]
Tìm số tự nhiên $y$, biết: $(y - 20\,000) : 4 = 15\,000$.
\shortans{80000}
\loigiai{
\textbf{Đáp số:} $80000$\\
$y - 20\,000 = 60\,000 \Rightarrow y = 80\,000$.
}
\end{ex}

\begin{ex} %[3D2V72]
Tính nhanh: $25\,000 \times 3 + 25\,000$.
\shortans{100000}
\loigiai{
\textbf{Đáp số:} $100000$\\
$25\,000 \times 4 = 100\,000$.
}
\end{ex}

\begin{ex} %[3D2V72]
Chia đều $84\,000\text{ kg}$ gạo cho 7 xã. Mỗi xã nhận được bao nhiêu ki-lô-gam gạo?
\shortans{12000}
\loigiai{
\textbf{Đáp số:} $12000$\\
$84\,000 : 7 = 12\,000\text{ (kg)}$.
}
\end{ex}

\begin{ex} %[3D2V72]
Một cửa hàng có $95\,000$ sản phẩm. Bán đi $\dfrac{1}{5}$ số sản phẩm đó. Hỏi cửa hàng còn lại bao nhiêu sản phẩm?
\shortans{76000}
\loigiai{
\textbf{Đáp số:} $76000$\\
Đã bán: $19\,000$. Còn lại: $95\,000 - 19\,000 = 76\,000$ sản phẩm.
}
\end{ex}
